% =====================================================================
%  Gaussian Random-Bridge Transport Numerical Experiments
%  Input this file in place of the [TO BE REWRITTEN] placeholder inside
%  \subsubsection{Gaussian Random-Bridge Transport Numerical Experiments}
%  Requires: booktabs, graphicx (both already loaded in main_Final.tex)
%  Requires the bibitems listed in bibitems_to_add.tex
% =====================================================================

We validate the training and simulation pipelines of Algorithms \ref{alg:training}
and \ref{alg:simulation} on noise-to-image generation over MNIST \cite{18} and
CIFAR-10 \cite{cifar2009}, and compare against simple denoising diffusion
\cite{12}, improved denoising diffusion \cite{13a}, their deterministic DDIM
samplers \cite{13b}, and conditional Flow Matching \cite{27}. Within each dataset every
model shares one network architecture, one exponential-moving-average decay, one
batch size and one training budget, and every sample is scored through one
Fr\'echet Inception Distance (FID) implementation \cite{fid2017}. Optimisation is
AdamW throughout at one learning rate and one schedule, with a single exception:
the random-bridge runs use zero weight decay while the three baselines use
AdamW's default $10^{-2}$, so the baselines are regularised slightly more
strongly than the bridge. Otherwise the differences between rows of the tables
below are the generative process and its simulation scheme.

\paragraph{Scope of the comparison.} Each row is a generative process simulated
directly from the dynamics it defines. Distillation-based approaches
\cite{19,21,22} operate in the same few-evaluation regime and reach lower FID
there, but they compress the sampling trajectory of an already-trained teacher
rather than defining a transport; they require that teacher and a second
training stage, and the same compression would apply to the random-bridge
construct. We therefore restrict the comparison to processes simulated directly
at a common training budget, and scale the empirical claims accordingly in
Section \ref{sec:gaussclaims}.

\paragraph{Configuration.} The network $f(\cdot\,;\boldsymbol{\Theta})$ is a UNet
with $64$ base channels on MNIST and $128$ on CIFAR-10, and self-attention at
the $16\times16$ and $8\times8$
resolutions, carrying $7.28$M parameters on MNIST and $52.54$M on CIFAR-10
(improved denoising diffusion carries $52.55$M, its learned reverse variance
doubling the output channels). Optimisation uses AdamW \cite{adamw2019} at
learning rate $2\times10^{-4}$ with a $500$-step linear warmup followed by a
constant rate, gradient-norm clipping at $1.0$, batch size $128$, $40\,000$
steps, an exponential moving average of the weights at decay $0.9999$, and
\texttt{bfloat16} autocast. Data is placed on a $32\times32$ grid in $[-1,1]$.
The terminal checkpoint is evaluated; no checkpoint is selected by FID. All runs
are executed on one NVIDIA GB10 accelerator with $119$GB of unified memory, on
which CIFAR-10 training proceeds at $2.84$ iterations per second, so that each
model costs $3.9$ hours to train.

\paragraph{Evaluation protocol.} Each table entry generates $50\,000$ samples and
compares their Inception-v3 pool3 statistics against those of the full training
split; FID against the test split is recorded alongside and moves in step with
it. Generated samples are mapped to 8-bit exactly once, by a clamped affine map,
and no per-image rescaling is applied to either real or generated images. The
matrix square root in the Fr\'echet distance is computed by eigendecomposition of
the symmetrised product; the implementation returns exactly $0$ on identical
inputs and $2048.0000$ on a unit analytic shift.

Because FID is biased upward at finite sample count, the achievable floor is
measured rather than assumed, using the comparator a table entry uses: $n$
held-out real images scored against the statistics of the full reference split.
On CIFAR-10 this gives $14.60$ at $n=2000$, $5.85$ at $n=5000$ and $3.15$ at
$n=10\,000$. The test split caps the held-out count at $10\,000$, so the floor at
$n=50\,000$ is extrapolated rather than measured: those three points follow
$20470\,n^{-0.955}$ to within a maximum log residual of $0.028$, placing the
$n=50\,000$ floor of Table \ref{tab:cifar} near $0.67$. Repeated simulation seeds give a standard deviation of $0.01$ to $0.07$
on the MNIST entries. On CIFAR-10 the seed spread is measured over three seeds at
$n=10\,000$, giving $0.12$ to $0.36$ at ten network evaluations and $0.05$ to
$0.09$ at fifty; since the sampling variance of FID decreases with $n$, a
standard deviation measured at $n=10\,000$ bounds the one at $n=50\,000$ from
above, and using it to decide which entries are resolved errs towards reporting
genuine differences as ties rather than the converse.

\paragraph{The regression target is the velocity field.}
Algorithm \ref{alg:training} regresses the $\mathcal{L}^2$ best-estimate
$\boldsymbol{Y}^{\boldsymbol{x}}_t$ under a loss that draws
$t\sim\mathcal{U}[0,T)$. Simulation, however, never consumes
$\boldsymbol{Y}^{\boldsymbol{x}}_t$ on its own: both Algorithm
\ref{alg:simulation} and the deterministic map \eqref{condflowmatch} advance the
state by a multiple of the velocity field
\begin{align}
\boldsymbol{u}(\xib^{\boldsymbol{x}}_t,t)
= \frac{\boldsymbol{Y}^{\boldsymbol{x}}_t - \xib^{\boldsymbol{x}}_t}{T-t},
\qquad t\in[0,T),
\label{eq:velocityrepeat}
\end{align}
already identified in \eqref{continuousvelocity}. An estimation error
$\boldsymbol{e}_t = f(\xib^{\boldsymbol{x}}_t,t;\boldsymbol{\Theta}) -
\boldsymbol{Y}^{\boldsymbol{x}}_t$ therefore enters the simulated trajectory at
magnitude $\|\boldsymbol{e}_t\|/(T-t)$, and a Girsanov comparison of the two
diffusions, which share the coefficient $\boldsymbol{\sigma}$, bounds the
Kullback-Leibler divergence of the generated terminal law from $\boldsymbol{\Psi}$
by $\int_0^T \E\|\boldsymbol{e}_t\|^2 (T-t)^{-2}\dd t$ up to a constant. The
objective of \eqref{optimizationlossfunction} minimises
$\int_0^T\E\|\boldsymbol{e}_t\|^2\dd t$. The two differ by the weight
$(T-t)^{-2}$ alone, and both are minimised at $\boldsymbol{e}_t\equiv\boldsymbol{0}$,
so the uniform objective identifies the correct function and the distinction is
immaterial at infinite capacity. At finite capacity it determines where the
unavoidable residual error is placed: under the uniform weight the terminal
region $t>0.98\,T$ contributes two per cent of the loss, and under $(T-t)^{-2}$
it contributes almost all of it.

Appendix \ref{app:weighting} measures the consequence. Two networks trained
identically, one against $\boldsymbol{Y}^{\boldsymbol{x}}_t$ and one against
$\boldsymbol{u}$, are the same estimator to three decimal places over the first
eighty per cent of the bridge and separate by a factor of ten at $t = 0.995\,T$;
in the units of \eqref{eq:velocityrepeat} the endpoint-trained error there is
$11.15$ against $0.70$. All results below therefore regress $\boldsymbol{u}$
against its conditional target
\begin{align}
\frac{\boldsymbol{y} - \xib^{\boldsymbol{x}}_t}{T-t}
= \frac{\boldsymbol{y}-\boldsymbol{x}}{T}
- \boldsymbol{\sigma}\sqrt{\frac{t}{T(T-t)}}\;\boldsymbol{\varepsilon},
\qquad \boldsymbol{\varepsilon}\sim\mathcal{N}(\boldsymbol{0},\boldsymbol{I}),
\label{eq:velocitytarget}
\end{align}
whose conditional mean given $\xib^{\boldsymbol{x}}_t$ is
$\boldsymbol{u}(\xib^{\boldsymbol{x}}_t,t)$ by the tower property, so the
regression is unbiased for the velocity field. The second term of
\eqref{eq:velocitytarget} has conditional variance growing like
$\boldsymbol{\sigma}^2 t/(T-t)$, whose integral diverges logarithmically;
training accordingly draws $t\sim\mathcal{U}[0,T-\epsilon)$ with
$\epsilon = 10^{-3}T$, mirroring the boundary offset already present in
Algorithm \ref{alg:simulation}. At $\boldsymbol{\sigma}=\boldsymbol{0}$ the
target reduces to $(\boldsymbol{y}-\boldsymbol{x})/T$ and the objective becomes
conditional Flow Matching exactly, consistent with Remark \ref{remflowmatch}:
the two schemes agree to $7\times10^{-7}$ in the simulated trajectory and to
machine precision in the loss. Flow Matching is therefore not an external
baseline in what follows but the $\boldsymbol{\sigma}\to\boldsymbol{0}$ member of
the same family, evaluated through the same code.

\paragraph{The diffusion scale and the horizon are one parameter.} Writing
$t = \tau T$ with $\tau\in[0,1]$, the conditional law of Corollary
\ref{exbrowniananticipative} reads
\begin{align}
\xib^{\boldsymbol{x}}_{\tau T} \mid \{\boldsymbol{Y}=\boldsymbol{y}\}
\;\sim\; \mathcal{N}\!\left( (1-\tau)\boldsymbol{x} + \tau\boldsymbol{y},\;
\boldsymbol{\sigma}^2 T\,\tau(1-\tau)\,\right), \notag
\end{align}
so training and simulation depend on $\boldsymbol{\sigma}$ and $T$ only through
the single scale $\boldsymbol{\gamma}^2 = \boldsymbol{\sigma}^2 T$. The
configuration $\boldsymbol{\sigma}=\boldsymbol{I}$, $T=0.1$ used below
corresponds to $\boldsymbol{\gamma} = 0.3162\,\boldsymbol{I}$, and the
sensitivity study of Appendix \ref{app:gamma} is conducted over
$\boldsymbol{\gamma}$ rather than over the two parameters separately.

\paragraph{Reporting convention.} Simulation cost is reported as the number of
network evaluations per generated sample, which equals the grid size $m$ of
Algorithm \ref{alg:simulation} for every scheme compared here. This is the
quantity that determines wall-clock time: all methods share one network and
spend one evaluation per step, so measured throughput at matched evaluation
count agrees to within one image per second (Section \ref{sec:gausscost}).

\subsubsection*{CIFAR-10}

\begin{table}[htbp]
\begin{center}
\caption{CIFAR-10. FID against the training split, $50\,000$ generated samples
per entry, extrapolated floor $0.67$ (measured $3.15$ at $n=10\,000$).
Simulation cost is the number of network
evaluations per sample. Bold marks the best entry in each column; a dash marks a
configuration not run.}
\label{tab:cifar}
\vskip 0.1in
\resizebox{\textwidth}{!}{\begin{tabular}{lccccccccc}
\toprule
Model & \multicolumn{9}{c}{Network evaluations per sample} \\
\cmidrule(lr){2-10}
 & 1 & 2 & 5 & 10 & 20 & 50 & 100 & 200 & 1000 \\
\midrule
Random-bridge, Euler--Maruyama & 385.98 & \textbf{168.38} & 43.99 & \textbf{19.21} & 12.26 & 9.83 & 9.17 & 8.88 & 9.23 \\
Flow Matching & 372.21 & 173.63 & \textbf{41.64} & 19.41 & 13.50 & 10.43 & 9.35 & 8.82 & 8.43 \\
Improved denoising, ancestral & \textbf{364.62} & 364.42 & 46.57 & 25.60 & \textbf{11.04} & \textbf{6.04} & \textbf{5.30} & \textbf{5.28} & \textbf{5.60} \\
Improved denoising, DDIM & 364.62 & 363.91 & 70.17 & 33.69 & -- & -- & -- & -- & -- \\
\bottomrule
\end{tabular}
}
\end{center}
\end{table}

Table \ref{tab:cifar} replicates the protocol throughout: the random-bridge row
uses $\boldsymbol{\sigma}=\boldsymbol{I}$ and $T=0.1$, simulated by Algorithm
\ref{alg:simulation} with the Euler--Maruyama scheme. Improved denoising
diffusion is reported with its ancestral sampler, which on CIFAR-10 is the
stronger of its two from five network evaluations upward; at one and two
evaluations the two coincide to within $0.6$ FID, both far from convergence.

From two to ten network evaluations the random-bridge transport attains lower
FID than improved denoising diffusion: $19.21$ against $25.60$ at ten
evaluations, and $168.38$ against $364.42$ at two, where ancestral sampling has
too few steps to depart from its prior. At a single evaluation neither has
departed from its prior and improved denoising is marginally ahead, $364.62$
against $385.98$. The ordering reverses between ten and twenty evaluations and
improved denoising leads thereafter, reaching $6.04$ against
$9.83$ at fifty and $5.28$ against $8.88$ at two hundred.

Against Flow Matching the two constructs track one another closely, as the
identity established above predicts. The bridge leads at two, ten, twenty, fifty
and one hundred evaluations and trails at one, five, two hundred and one
thousand; from ten evaluations upward the largest difference in either
direction is $1.24$ FID, at twenty. At the one- and two-evaluation extremes,
where neither construct has converged, the differences are larger: $13.77$ in
favour of Flow Matching at one and $5.25$ in favour of the bridge at two.
At ten evaluations the gap of $0.20$ lies inside one standard
deviation of the simulation-seed spread and is reported as a tie, whereas at
fifty the gap of $0.60$ is roughly seven standard deviations and is resolved.

Generation quality improves monotonically through two hundred evaluations,
where it attains $8.88$, and rises by $3.9$ per cent to $9.23$ at one thousand.
This terminal rise is not specific to the random-bridge construct: improved
denoising diffusion rises by $5.9$ per cent over the same range, from $5.28$ to
$5.60$, while Flow Matching, which integrates a deterministic flow, decreases
throughout and attains $8.43$. The effect is therefore attributable to
stochastic simulation at fine discretisation rather than to the transport.

\subsubsection*{MNIST}

\begin{table}[htbp]
\begin{center}
\caption{MNIST. FID against the training split, $50\,000$ generated samples per
entry; the measured floor is $0.98$ at $n=10\,000$ and does not admit a
defensible extrapolation to $50\,000$. Mean $\pm$ standard deviation over
simulation seeds where repeats were run.}
\label{tab:mnist}
\vskip 0.1in
\resizebox{\textwidth}{!}{\begin{tabular}{lccccccccc}
\toprule
Model & \multicolumn{9}{c}{Network evaluations per sample} \\
\cmidrule(lr){2-10}
 & 1 & 2 & 5 & 10 & 20 & 50 & 100 & 200 & 1000 \\
\midrule
Random-bridge, exact transition & 346.04 & 50.31 & 9.62 & 4.08 $\pm$ 0.06 & 2.14 & 1.37 $\pm$ 0.00 & 1.19 & 1.12 $\pm$ 0.00 & -- \\
Random-bridge, Euler--Maruyama & 346.30 & \textbf{49.42} & \textbf{9.46} & \textbf{3.72 $\pm$ 0.02} & \textbf{1.89} & 1.27 & 1.17 & 1.12 & 1.06 \\
Random-bridge, deterministic map & 346.05 & 51.46 & 10.51 & 5.73 $\pm$ 0.08 & 4.19 & 3.42 & 3.19 & 3.05 & -- \\
Flow Matching & 298.68 & 54.09 & 10.34 & 4.13 $\pm$ 0.06 & 2.05 & \textbf{1.23 $\pm$ 0.02} & \textbf{1.07} & \textbf{1.02 $\pm$ 0.01} & \textbf{1.02} \\
Simple denoising, cosine, DDIM & 418.50 & 380.07 & 17.76 & 6.96 & 4.07 & 2.24 & 1.40 & 2.83 & 16.30 \\
Simple denoising, cosine, ancestral & 377.15 & 380.14 & 16.27 & 9.01 & 4.83 & 3.00 & 2.90 & 2.99 & 2.71 \\
Simple denoising, linear, DDIM & 418.36 & 285.10 & 14.05 & 7.22 & 4.48 & 2.79 & 1.67 & 2.74 & 13.26 \\
Simple denoising, linear, ancestral & \textbf{282.19} & 284.25 & 18.14 & 10.85 & 6.87 & 3.94 & 3.19 & 3.07 & 3.05 \\
Improved denoising, ancestral & 389.46 & 393.23 & 20.06 & 10.99 & 7.16 & 3.62 & 2.02 & 1.74 & 1.93 \\
Improved denoising, DDIM & 389.46 & 392.03 & 23.62 & 7.82 & 4.22 & 2.38 & 1.24 & 2.72 & 21.94 \\
\bottomrule
\end{tabular}
}
\end{center}
\end{table}

On MNIST the random-bridge transport attains lower FID than improved denoising
diffusion at every evaluation count: $3.72$ against $10.99$ at ten evaluations
and $1.06$ against $1.93$ at one thousand. Taken across both of its samplers the
strongest baseline entry is $1.24$, obtained by improved denoising with DDIM at
one hundred evaluations, against $1.06$ for the bridge. Against Flow Matching the
bridge leads from two to twenty evaluations ($49.42$ against $54.09$ at two;
$9.46$ against $10.34$ at five; $3.72$ against $4.13$ at ten; $1.89$ against
$2.05$ at twenty), trails at a single evaluation ($346.30$ against $298.68$),
and trails from fifty upward ($1.27$ against $1.23$; $1.17$ against
$1.07$; $1.12$ against $1.02$; $1.06$ against $1.02$). The ordering with respect
to improved denoising differs from CIFAR-10, where the bridge leads only from
two to ten evaluations.

Both DDIM baselines deteriorate past one hundred evaluations, reaching $16.30$
and $13.26$ at one thousand. Clipping the forward variance at $0.999$ leaves the
cosine schedule with a terminal signal-to-noise ratio of
$\bar\alpha_T = 2.4\times10^{-9}$, at which noise prediction is degenerate and
the reconstruction of the clean sample divides by
$\sqrt{\bar\alpha_T} = 4.9\times10^{-5}$ \cite{snr2024}. The effect is confined
to the deterministic sampler: the ancestral samplers are flat over the same
range, and two independent timestep grids reproduce it. The prescribed remedy is
velocity prediction, which we do not apply; the rows are reported as measured
because a baseline that deteriorates would otherwise inflate the margin over it.

\subsubsection*{Fidelity and coverage}
\label{sec:gausscoverage}

FID compresses sample fidelity and distributional coverage into a single scalar,
so a model producing clean samples of part of $\boldsymbol{\Psi}$ and one
producing diffuse samples of all of it may score alike. Appendix
\ref{app:coverage} separates the two using precision and recall
\cite{prec2019} together with the outlier-robust density and coverage
\cite{dens2020}, computed from the same Inception activations, and reports the
Kernel Inception Distance \cite{binkowski2018}, which is unbiased and therefore not
subject to the finite-sample bias that sets the FID floor.

Three findings follow. First, additional network evaluations purchase coverage
rather than fidelity: for the random-bridge transport, precision is flat across
the range at $0.606$, $0.614$ and $0.607$ for ten, fifty and two hundred
evaluations, while coverage rises from $0.571$ to $0.638$ to $0.680$. Flow
Matching and improved denoising diffusion exhibit the same pattern. Second,
where the bridge leads at low evaluation count it leads on recall: at ten
evaluations it is level with improved denoising on precision ($0.606$ against
$0.617$) and on coverage ($0.571$ against $0.578$), and the FID advantage is
carried by recall ($0.530$ against $0.490$). Third, the Kernel Inception
Distance orders every row identically to FID at every evaluation count, so the
ordering is not an artefact of the FID estimator.

Because the Inception network was fitted to natural colour photographs, the
MNIST comparison is repeated in a feature space fitted to the data: a
convolutional classifier trained on the MNIST training split to $97.8$ per cent
test accuracy. The ordering is unchanged, and the classifier's predicted-label
histogram shows that no model omits a digit, every model serving its least
frequent digit at $92$ per cent or more of the real rate at one hundred
evaluations (Appendix \ref{app:coverage}).

\subsubsection*{Simulation cost at matched quality}
\label{sec:gausscost}

\begin{table}[htbp]
\begin{center}
\caption{Network evaluations per generated sample required to reach a given FID
on CIFAR-10, by logarithmic interpolation between the measured entries of Table
\ref{tab:cifar}.}
\label{tab:cost}
\vskip 0.1in
\begin{tabular}{lccc}
\toprule
Target FID & Random-bridge & Flow Matching & Improved denoising \\
\midrule
$30$ & $7$   & $7$   & $8$  \\
$20$ & $10$  & $10$  & $12$ \\
$15$ & $15$  & $16$  & $16$ \\
$12$ & $22$  & $30$  & $19$ \\
$10$ & $47$  & $65$  & $23$ \\
$9$  & $151$ & $157$ & $27$ \\
\bottomrule
\end{tabular}
\end{center}
\end{table}

Since every scheme compared here evaluates the same network once per step,
throughput at matched evaluation count is equal to within measurement noise:
$176$, $177$ and $177$ images per second at ten evaluations for the
random-bridge, Flow Matching and improved denoising respectively. Peak allocated
memory likewise agrees, at $8645$, $8647$ and $8648$ MB for one training step at
batch $128$ and $3521$, $3510$ and $3545$ MB for one simulation batch of $500$ at
fifty evaluations. Evaluation count therefore converts proportionally into
wall-clock time, and the discriminating comparison is cost at equal quality.

Table \ref{tab:cost} reports it. The random-bridge construct is the least
expensive route to FID $15$ and above and the most expensive to FID $10$ and
below. At FID $20$ the advantage over improved denoising is a factor of $1.2$ in
network evaluations.

\subsubsection*{Summary of empirical findings}
\label{sec:gaussclaims}

\begin{enumerate}
\item On CIFAR-10, at matched architecture, training budget and evaluation code,
the random-bridge transport attains lower FID than improved denoising diffusion
from two to ten network evaluations ($19.21$ against $25.60$ at ten) and higher
FID from twenty upward ($9.83$ against $6.04$ at fifty). On MNIST it attains lower FID
at every evaluation count tested ($1.06$ against $1.93$ at one thousand).
\item Against Flow Matching the two are close throughout. On MNIST the bridge
leads from two to twenty evaluations and trails at one and from fifty upward. On CIFAR-10
it leads at two and from ten to one hundred, and trails at one, five, two
hundred and one thousand; from ten evaluations upward the largest difference in
either direction is $1.24$ FID.
\item Quality improves monotonically through two hundred evaluations and rises
by $3.9$ per cent at one thousand, a terminal rise improved denoising diffusion
also exhibits at $5.9$ per cent and deterministic Flow Matching does not.
\item Additional evaluations purchase coverage rather than fidelity, and the
advantage at low evaluation count is carried by recall.
\item Cost at equal quality favours the random-bridge construct by a factor of
$1.2$ in network evaluations at FID $20$ and favours improved denoising by a
factor of $2$ at FID $10$.
\end{enumerate}

These statements concern generative processes simulated directly, at one shared
training budget, and support a favourable trade-off between simulation cost and
generation quality in the few-evaluation regime. They do not support a claim of
best absolute FID, nor of an order-of-magnitude efficiency advantage.
